扩散模型生成一张图片，做的事情说穿了是解一个随机微分方程：从一团纯高斯噪声出发，沿着某个学来的向量场一步一步走，走到终点就是一张图。这句话里几乎每个词都欠着交代——随机微分方程是什么意思，一团噪声怎么会知道往哪边走，那个向量场是什么东西，它凭什么能从数据里学出来。这一章把这句话拆开。

前面几章的 Markov 链在离散状态之间跳：离散时间靠转移矩阵 \(P\)，连续时间靠生成矩阵 \(Q\)，演化由 \(P^n\) 或 \(e^{tQ}\) 推进；时间序列虽然取连续值，讨论的仍是一个个离散时刻。现在时间和状态都连续。位置、价格、浓度、参数向量的变化来自无数微小扰动的持续累积，而不是一次一次的跳跃。有限矩阵让位给微分算子，\(e^{tQ}\) 的位置被偏微分方程和随机微分方程生成的演化接手。

这个过渡比``把求和换成积分''难得多。极限过程本身带着离散世界没有的粗糙性：轨道处处连续，却处处不可微。粗糙性不是需要绕开的技术麻烦，它恰恰是这一章所有反常公式的出处，所以第一步是把它看清楚。

先用一个微型算例感受尺度。无偏 \(\pm1\) 随机游走走 400 步，位置方差是 400，标准差 20。若整段过程要压进 1 秒，即 \(\Delta t=1/400\)，那么步长必须同步缩到 \(\sqrt{1/400}=1/20\)，重跑一遍位置方差恰好是 1，得到标准正态。步长缩得再快一点，方差蒸发成零；缩得慢一点，方差炸开。留给极限的只有 \((\Delta x)^2/\Delta t\) 保持有限这一条缝——这条缝就是整章的地基。

各篇沿两条互补的路线展开。路径路线是随机游走走向布朗运动、再走向 Itô 积分与随机微分方程；分布路线是局部的漂移与扩散走向 Fokker--Planck 方程、再走向得分与反向生成。

\begin{itemize}
\tightlist
\item \hyperref[sec:06-Random-Processes/06-Diffusion-and-SDE/01]{``从随机游走到布朗运动''} 建立 \(\Delta x\sim\sqrt{\Delta t}\) 这条标度关系，并说明为什么它同时导致路径不可微和二次变差非零；
\item \hyperref[sec:06-Random-Processes/06-Diffusion-and-SDE/02]{``Fokker--Planck 方程描述密度的流动''} 换到密度视角，单条轨道随机而密度的演化是确定的偏微分方程，顺带交出 Langevin 采样与温度的角色；
\item \hyperref[sec:06-Random-Processes/06-Diffusion-and-SDE/03]{``Itô 积分与随机微分方程''} 重建被粗糙性废掉的积分与链式法则，为此多出一个二阶修正项；
\item \hyperref[sec:06-Random-Processes/06-Diffusion-and-SDE/04]{``数值模拟与时间反演的直觉''} 给出 Euler--Maruyama 格式与两种收敛，再把时间倒过来，让得分 \(\nabla\log p_t\) 登场；
\item \hyperref[sec:06-Random-Processes/06-Diffusion-and-SDE/05]{``得分匹配：怎样从样本学习得分''} 补上最后一块：得分能从样本估出来，而且绕开了归一化常数。
\end{itemize}

第一篇的标度关系是全章的枢纽，值得多花时间；后面每次遇到``为什么这里会多出一项''，答案基本都能追回到它。若目标只是搞懂扩散模型的机制，Itô 积分的构造细节可以先跳过，但时间方向、扩散系数的形式、得分对应的是哪一个边缘分布这三件事不能含糊，它们是实现里最常出错的地方。

严格性的边界也先说明白。Itô 积分的完整构造、SDE 解的存在唯一性、时间反演定理，都依赖测度论，超出本书范围。这一章的交付方式是分层的：直觉与构造思想完整给出，正式定理声明结论与成立条件但不展开证明，并标出边界在哪里。读完之后你应当能回答布朗运动为什么那样定义、Fokker--Planck 每一项在搬运什么、Itô 修正项从哪里来、扩散模型为什么能从噪声走回数据，以及得分怎样从样本中学到；至于逐条证明它们为什么成立，留给更严格的教材。
